\expandafter\def\csname HistoricalNoteText363\endcsname{Dieser Aufsatz enthält den wesentlichen Inhalt des Vortrags, den ich am 19. September 1837 in der öffentlichen Sitzung der Königlichen Sozietät der Wissenschaften gehalten habe.}
\clearpage
\fancyhead[L]{\small ÜBER EIN NEUES INSTRUMENT.}
\SourceLeaf{363}
\pdfbookmark[0]{Über ein neues Instrument. 1837}{instrument-1837}
\begin{center}
{\large ÜBER EIN NEUES INSTRUMENT, DAS IN ERSTER LINIE ZUR UNMITTELBAREN BEOBACHTUNG DER ÄNDERUNGEN IN DER INTENSITÄT DES HORIZONTALEN ANTEILS DES ERDMAGNETISMUS BESTIMMT IST\HistoricalNote{363}.}
\par\medskip
Resultate aus den Beobachtungen des magnetischen Vereins. 1837. I.
\end{center}

Zur vollständigen Bestimmung des Erdmagnetismus an einem gegebenen Ort braucht man bekanntlich ein System von \textit{drei} Elementen. Gewöhnlich wählt man dafür die Abweichung, die Neigung und die Stärke. Diese Wahl ist zwar begrifflich die einfachste, doch lässt sich auch eine andere Kombination zugrunde legen, die in mancher Hinsicht sogar vorzuziehen sein kann. Insbesondere ist es sowohl praktisch als auch theoretisch weit vorteilhafter, den horizontalen Anteil der erdmagnetischen Kraft für sich zu betrachten und durch zwei Elemente darzustellen, nämlich durch seine Richtung (Deklination) und seine Stärke. Nimmt man dann entweder die Stärke der vertikalen Kraft oder die Neigung der gesamten Kraft als drittes Element hinzu, ergibt sich daraus bei Bedarf unmittelbar die Stärke der gesamten Kraft.

Was nun die beiden Elemente des horizontalen Erdmagnetismus betrifft, um den es hier allein geht, so hat das seit fünf Jahren eingesetzte Magnetometer alle Aufgaben, die sich bei der Deklination stellen, voll\SourceLeaf{364}ständig gelöst. Es dient zur Bestimmung ihres absoluten Werts ebenso wie zur Verfolgung ihrer regelmäßigen und zufälligen Änderungen von Jahr zu Jahr, von Monat zu Monat, von Stunde zu Stunde, ja selbst von einer Minute zur nächsten, und zwar mit einer Sicherheit, Bequemlichkeit und Genauigkeit, die nichts zu wünschen übrig lassen.

Dasselbe Instrument dient zugleich dazu, die Stärke des horizontalen Erdmagnetismus in absolutem Maß zu bestimmen. Gerade diese Aufgabe hat bekanntlich überhaupt erst zur Einrichtung des Magnetometers Anlass gegeben. Gleichwohl löst es die Aufgabe bei Weitem noch nicht in \textit{jeder} Hinsicht vollständig.

Um deutlich zu machen, was dabei noch zu wünschen bleibt, muss ich zunächst daran erinnern, dass die Anwendung des Magnetometers zur Bestimmung der magnetischen Intensität auf einer Verbindung \textit{mehrerer} Arbeitsschritte beruht. Einer davon besteht darin, die Schwingungsdauer einer Nadel zu beobachten. Diese Beobachtung nimmt ihrer Natur nach eine beträchtliche Zeit in Anspruch, weil die Zahl der Schwingungen, aus denen man auf die Dauer einer einzelnen schließen muss, nicht zu klein sein darf. Bleibt die magnetische Intensität während eines solchen Vorgangs konstant, so entspricht die gefundene Schwingungsdauer tatsächlich \textit{diesem} Wert der Intensität. Hat sie sich inzwischen geändert, entspricht die Schwingungsdauer nur \textit{dem Mittelwert} der Intensität in diesem Zeitraum. Auf diese Weise bleibt uns aber völlig verborgen, ob sich die magnetische Intensität \textit{während} dieser Zeit überhaupt geändert hat und welche Änderungen aufgetreten sind. Das Instrument kann also nur Durchschnittswerte über bestimmte Zeiträume liefern, ohne den vollständigen Verlauf innerhalb dieser Zeiträume getreu wiederzugeben. Wollte man diesem Verlauf näherkommen, indem man kürzere Zeiträume wählte oder die Ergebnisse immer nur auf wenige Schwingungen stützte, so würden die Ergebnisse zu sehr an Genauigkeit und Sicherheit verlieren. Man liefe Gefahr, bloße Beobachtungsfehler für Anomalien der Intensität zu halten.

Je interessanter sich nun gerade die in kurzen Zeitabständen wechselnden Störungen der erdmagnetischen Kraft nach den Erfahrungen der letzten Jahre schon in ihrer einseitigen Erscheinung bei der Deklination gezeigt haben, desto dringender wünschte man sich ein Mittel, mit dem sich auch die nicht \SourceLeaf{365}zu bezweifelnden Wirkungen solcher Störungen auf die Intensität mit derselben Leichtigkeit, Sicherheit und Genauigkeit verfolgen und messen ließen.,

Wie ich eben dargelegt habe, eignen sich die bisherigen Beobachtungsmittel für \textit{diesen} Zweck deshalb nicht, weil sie auf Beobachtungen von Schwingungszeiten beruhen. Solche Beobachtungen dauern ihrer Natur nach jedes Mal zu lange. Die Schwingungsdauer einer Nadel dient hier aber selbst nur dazu, mittelbar das Drehmoment zu bestimmen, das die erdmagnetische Kraft auf die Nadel ausübt, wenn sie sich außerhalb des magnetischen Meridians befindet. Lässt sich dieses Drehmoment also auf direktem Weg, ohne Schwingungen zu beobachten, genau bestimmen und lassen sich seine Änderungen sicher, genau und schnell messen, dann ist unsere Aufgabe im Wesentlichen gelöst. Das Mittel, das ich dafür verwendet habe, beruht auf folgender Grundlage.

Die Bedingungen für das Gleichgewicht eines an zwei Fäden aufgehängten Körpers beliebiger Gestalt lassen sich kurz zusammenfassen, wenn man vorerst annimmt, dass seine Teile nur der Schwerkraft unterliegen und fest miteinander verbunden sind. Die Vertikale durch den Schwerpunkt des Körpers und die durch die Fäden dargestellten Geraden müssen in einer Ebene liegen und zugleich entweder zueinander parallel sein oder sich in einem Punkt schneiden. In der Gleichgewichtsstellung liegen die beiden Fäden und der Schwerpunkt also immer in einer vertikalen Ebene. Um sich das klar vorzustellen, kann man annehmen, dass die beiden Fäden gleich lang sind, ihre oberen Befestigungspunkte auf gleicher Höhe liegen und ebenso weit voneinander entfernt sind wie die beiden unteren, und schließlich, dass die beiden unteren Punkte mit dem Schwerpunkt ein gleichschenkliges Dreieck bilden. Unter diesen Voraussetzungen hängen die beiden Fäden im Gleichgewicht vertikal, und eine dritte, genau in ihrer Mitte gedachte Vertikale geht durch den Schwerpunkt des Körpers. Dreht man den Körper um diese letzte Gerade aus seiner Lage, verlaufen die beiden Fäden weder vertikal noch in einer Ebene, und zugleich wird der Körper etwas angehoben. Dadurch entsteht das Bestreben, in die vorige Lage zurückzukehren, mit einem Drehmoment, das man mit hinreichender Genauigkeit als proportional zum Sinus der Auslenkung aus der Ruhelage ansetzen kann. Es ist also am größten, wenn die Auslenkung 90 Grad beträgt. Dieses größte Drehmoment ist im\SourceLeaf{366}mer stillschweigend gemeint, wenn ohne weitere Bestimmung von einem Drehmoment gesprochen wird. Man kann es auch als Maß einer Kraft ansehen, mit der der Körper aufgrund seiner Aufhängung in seiner Gleichgewichtsstellung gehalten wird. Der Kürze halber will ich sie die durch die Aufhängung entstehende Richtkraft nennen. Ihre Größe hängt 1) von der Länge der Aufhängungsfäden, 2) von ihrem Abstand, 3) vom Gewicht des Körpers ab. Sie ist umgekehrt proportional zur Länge der Fäden, direkt proportional zum Quadrat ihres Abstands und ebenfalls direkt proportional zum Gewicht des Körpers. Treffen die obigen Voraussetzungen nicht genau zu, ist der Ausdruck für die Richtkraft komplizierter. Auch die Reaktion der Fäden auf eine Verdrillung macht noch eine kleine Abänderung nötig. Es gibt jedoch Möglichkeiten, die Größe der Richtkraft durch Versuche mit größter Genauigkeit zu bestimmen. Überlässt man den Körper nach einer kleineren oder größeren Auslenkung aus der Gleichgewichtsstellung sich selbst, so schwingt er mit größter Regelmäßigkeit. Die Mitte seiner Schwingungen fällt mit dieser Stellung zusammen, und ihre Dauer hängt von der Größe der Richtkraft und vom Trägheitsmoment des Körpers ab.

Nehmen wir nun an, dass ein horizontaler Magnetstab einen Teil des aufgehängten Körpers bildet, so kommt eine zweite Richtkraft ins Spiel. Die Erscheinungen hängen dann nach den bekannten Regeln der Statik von der Zusammensetzung der beiden Richtkräfte ab. Dabei sind drei Fälle zu unterscheiden. Die beiden Stellungen, in denen der Körper unter der Wirkung jeweils einer der beiden Kräfte allein im Gleichgewicht wäre, können zusammenfallen, einander entgegengesetzt sein oder einen Winkel miteinander bilden. Man sieht leicht, dass der Unterschied zwischen diesen drei Fällen vom Verhältnis zweier Winkel abhängt. Den einen bildet die Gerade durch die beiden unteren Befestigungspunkte der Fäden mit dem Magnetstab, den anderen die Gerade durch die beiden oberen Befestigungspunkte mit dem magnetischen Meridian. Denkt man sich den Körper in der allein durch die Aufhängung bestimmten Gleichgewichtslage, so muss der Magnetstab im ersten unserer drei Fälle im magnetischen Meridian in seiner natürlichen Lage liegen (der Nordpol auf der Nordseite). Im zweiten Fall muss er in umgekehrter Lage im Meridian liegen und im dritten einen Winkel mit dem magne\SourceLeaf{367}tischen Meridian bilden. Der Kürze halber will ich diese drei möglichen Lagen des Magnetstabs im Apparat die natürliche, die umgekehrte und die transversale nennen.

In der natürlichen Lage ändert die Wirkung des Erdmagnetismus auf den Magnetstab die der Aufhängung entsprechende Gleichgewichtsstellung des Apparats nicht. Der Apparat wird aber mit einer verstärkten Kraft darin gehalten, die der Summe der beiden Richtkräfte entspricht.

Im zweiten Fall, der umgekehrten Lage, bleibt das Gleichgewicht in dieser Stellung ebenfalls bestehen. Es ist aber nur stabil, wenn die magnetische Richtkraft kleiner ist als die Richtkraft der Aufhängung. Der Apparat wird dann nur mit einer Kraft in dieser Stellung gehalten, die der Differenz der beiden Richtkräfte entspricht. Wäre dagegen die magnetische Richtkraft größer, so wäre das Gleichgewicht instabil. Einmal ausgelenkt, würde der Apparat nicht dorthin zurückkehren. Er würde sich immer weiter von dieser Stellung entfernen und erst in der entgegengesetzten Stellung zur Ruhe kommen, in der der Magnetstab seine natürliche Lage im Raum hat, die Aufhängungsfäden einander aber kreuzen.

Im dritten Fall schließlich, in dem die beiden Richtkräfte einen Winkel miteinander bilden, kommt es zwischen ihnen zu einem Ausgleich in einer Zwischenstellung. Dabei liegt weder der Magnetstab im Meridian noch verläuft die Gerade durch die unteren Befestigungspunkte der Fäden parallel zu der durch die oberen. Sowohl diese Zwischenlage als auch die Kraft, mit der der Apparat darin gehalten wird, richten sich nach dem Gesetz der Statik für die Zusammensetzung zweier Kräfte. Nun ist zugleich ersichtlich, dass sich das Verhältnis der beiden beteiligten Richtkräfte berechnen lässt, wenn der Apparat eine Möglichkeit bietet, die Winkel zwischen den drei fraglichen Stellungen zu messen. Folglich kann man die magnetische Richtkraft auch in absolutem Maß angeben, wenn die Richtkraft der Aufhängung in absolutem Maß bekannt ist. Damit ist unsere Aufgabe gelöst. Am vorteilhaftesten ist es übrigens, den Magnetstab im Verhältnis zu den anderen Teilen des Apparats so einzulegen, dass er in der so entstehenden Gleichgewichtsstellung beinahe einen rechten Winkel mit dem magnetischen Meridian bildet. Für diesen Fall ist die Bezeichnung transversale Lage besonders treffend. Zum einen ist dann die Abweichung der Fäden \SourceLeaf{368}von ihrer Lage in einer Ebene am größten, sodass sich das Ergebnis mit der größten Genauigkeit berechnen lässt. Zum anderen hat dann die kleine Änderung der magnetischen Deklination infolge der stündlichen oder zufälligen Schwankungen keinen merklichen Einfluss auf die Stellung. Dagegen beeinflusst jede Änderung der Stärke des Erdmagnetismus die Stellung unmittelbar und lässt sich sogleich mit derselben Leichtigkeit, Schnelligkeit und Genauigkeit erkennen und messen wie das Wechselspiel der Deklinationsänderungen am gewöhnlichen Magnetometer.

Die praktische Anwendbarkeit dieser Idee hatte ich schon vor mehreren Jahren durch vorläufige Versuche mit einer allerdings noch sehr groben Vorrichtung bestätigt gefunden. Darauf habe ich auch in meinem Aufsatz über Erdmagnetismus und Magnetometer hingewiesen [S. 327 dieses Bandes]. Vor Kurzem habe ich jedoch einen vollkommeneren Apparat anfertigen lassen und ihn in der Sternwarte an der Stelle aufgehängt, an der sich bisher das Magnetometer mit dem fünfundzwanzig Pfund schweren Stab befand. Nach den bisherigen Erläuterungen lässt sich dieser Apparat kurz beschreiben.

Er hängt an zwei 17 Fuß langen Stahldrähten, oder genauer gesagt an einem einzigen Draht, dessen Enden unten am Apparat befestigt sind. Seine Mitte läuft oben über zwei Zylinder, die seine beiden Stränge in einem geeigneten Abstand halten (etwa $1\frac{1}{2}$ Zoll). Diese Einrichtung hat zugleich den Vorteil, dass die beiden Stränge von selbst die gleiche Spannung haben. Die Aufhängung liegt oberhalb der Saaldecke, und die Drähte hängen frei durch eine kreisrunde Öffnung von $3\frac{1}{2}$ Zoll Weite in der Decke. Der Abstand zwischen den Drähten kann oben wie unten nach Bedarf vergrößert oder verkleinert werden. Der an den Drähten hängende Apparat selbst besteht aus vier Hauptteilen. Der erste, an dem die Drähte befestigt sind, ist eine horizontale Kreisscheibe von vier Zoll Durchmesser mit einer Einteilung in Viertelgrade auf Silber. Der zweite Teil besteht aus einer Alhidade, die sich konzentrisch auf dem Teilkreisrand drehen lässt und zwei Nonien zum Ablesen einzelner Winkelminuten hat. Mit der Alhidade ist eine recht kräftige Stange fest verbunden, die senkrecht zur Kreisebene steht. An ihr sitzt ein sehr genau kreisrunder Spiegel von $1\frac{1}{2}$ Zoll Durchmesser. Darin sieht man durch ein 16 Fuß entferntes Fernrohr das Bild eines Ausschnitts einer horizontalen, in einzelne Millimeter eingeteilten Skala, die unter dem Fernrohr befestigt ist. Auf diese Weise lässt sich \SourceLeaf{369}also jede Änderung der Lage des Kreises erkennen und messen. Kleine Änderungen misst man unmittelbar und mit größter Genauigkeit an den im Fernrohr sichtbaren Skalenteilen, größere, indem man damit eine Bewegung der Alhidade verbindet und die Nonien abliest. Der dritte Teil ist das Schiffchen unter dem Kreis, ein doppelter Rahmen, durch den der vierte Bestandteil, ein starker, fünfundzwanzig Pfund schwerer Magnetstab, gesteckt wird. Auch dieses Schiffchen ist um den Mittelpunkt des Kreises drehbar. Es hat zwei auf dem Teilkreisrand aufliegende Nonien, mit denen man die Größe der Drehung auf die Winkelminute genau messen kann.

Stellt man nun zunächst das Schiffchen so ein, dass der Apparat dieselbe Gleichgewichtslage behält, ob im Schiffchen der Magnetstab oder ein nichtmagnetischer Körper vom selben Gewicht liegt, so erhält man die erste oder die zweite der zuvor unterschiedenen Hauptlagen, je nachdem, ob der Magnetstab dabei in seiner natürlichen oder in seiner umgekehrten Lage ist. Die erste bietet keine besonders wichtige praktische Anwendung. Die Brauchbarkeit der zweiten hängt davon ab, dass die magnetische Richtkraft etwas kleiner ist als die Richtkraft der Aufhängung. Beim hiesigen Apparat ist das Verhältnis dieser Richtkräfte jetzt nahezu 10 zu 11. Die resultierende Richtkraft beträgt also nur ein Zehntel der magnetischen Richtkraft. Wir haben hier somit ein Gegenstück zu einer astatischen Magnetnadel. Jede äußere Kraft, die die Richtung einer einfachen Nadel stört, hat hier eine zehnmal größere Wirkung als bei der Aufhängung an einem einzigen Faden, und zwar, wie man leicht einsieht, in entgegengesetzter Richtung. Damit ist unter anderem eine Aufgabe gelöst, mit der man sich früher wiederholt ohne Erfolg beschäftigt hat, nämlich die täglichen und stündlichen Schwankungen der magnetischen Deklination vergrößert darzustellen. Wiederholte gleichzeitige Beobachtungen dieser Art mit diesem Apparat und mit dem Magnetometer des magnetischen Observatoriums haben zwar immer höchst befriedigende Ergebnisse geliefert. Inzwischen verliert diese Anwendung jedoch an Bedeutung, weil die gewöhnlichen Magnetometer schon die kleinsten Änderungen mit aller gewünschten Genauigkeit anzeigen und daher eine Vergrößerung nicht mehr nötig ist.

Diese und andere Anwendungen bei umgekehrt eingelegtem Magnetstab, auf \SourceLeaf{370}die ich später noch zurückkommen werde, sind jedoch nur von untergeordneter Bedeutung. Weitaus wichtiger ist die Anwendung des Apparats in der dritten oder \textit{transversalen} Lage für Intensitätsmessungen. Bringt man den Magnetstab ausgehend von der natürlichen Lage durch Drehen des Schiffchens aus dem magnetischen Meridian, muss sich der ganze Apparat bis zum Gleichgewicht um einen bestimmten Winkel zurückdrehen, der dem Verhältnis der beiden Richtkräfte entspricht. Die Differenz der beiden Winkel ist dann die Abweichung des Magnetstabs vom magnetischen Meridian in der Gleichgewichtsstellung. Man kann es leicht so einrichten, dass diese Abweichung beinahe 90 Grad beträgt, und so die bereits genannten Vorteile erzielen. Der Apparat eignet sich dann ganz besonders zur Beobachtung der \textit{Änderungen} der Intensität, die sich unmittelbar an seiner veränderten Stellung zeigen. Bei Änderungen, die sich erst nach längerer Zeit einstellen, müssen der Natur der Sache nach mehrere Umstände berücksichtigt werden. Insbesondere muss von Zeit zu Zeit mit geeigneten (bekannten) Mitteln untersucht werden, ob und in welchem Maß sich die Stärke des Magnetismus im Stab verändert hat. Auch Temperaturänderungen sind zu berücksichtigen, sowohl soweit sie diese Stärke beeinflussen als auch soweit sie den Abstand und die Länge der Aufhängungsdrähte und damit die Richtkraft der Aufhängung beeinflussen. Bei den unregelmäßigen Änderungen der Intensität, die in kurzen Zeitabständen wechseln, leistet der Apparat jedoch genau dasselbe wie das Magnetometer bei ähnlichen Änderungen der Deklination. Auch die Beobachtungsweise ist bei beiden Apparaten genau gleich. Man erhält die Änderungen der Intensität zunächst in Skalenteilen, die sich leicht in Bruchteile der Intensität selbst umrechnen lassen. Unter den gegenwärtigen Verhältnissen des Apparats entspricht einem Skalenteil der 22000\textsuperscript{ste} Teil der gesamten Intensität.

Die bisherigen Erfahrungen mit dem Apparat umfassen zwar erst eine kurze Zeit und sind noch nicht sehr zahlreich, lassen aber schon einige durchaus wichtige Ergebnisse erkennen.

Erstens deuten die bisherigen Beobachtungen auf regelmäßige Änderungen hin, die von der Tageszeit abhängen. Sie mögen sich allerdings ebenso häufig mit unregelmäßigen Änderungen vermischen wie bei der Deklination, und ihre sichere Unterscheidung \SourceLeaf{371}wird jahrelang fortgesetzte Beobachtungen erfordern. Wenn ich bei den wenigen bisherigen Erfahrungen eher eine Vermutung als ein Ergebnis aussprechen darf, scheint der regelmäßige Verlauf darin zu bestehen, dass die Intensität im Laufe des Vormittags abnimmt, schon eine oder zwei Stunden vor Mittag ihr Minimum erreicht und von da an wieder zunimmt. Um wenigstens vorläufig einen Anhaltspunkt für das quantitative Verhältnis zu gewinnen, habe ich im August 1837 an 30 Tagen die Stellung morgens um 10 Uhr und nachmittags um 3 Uhr aufgezeichnet. An 26 Tagen war die Intensität nachmittags größer und nur an 4 Tagen kleiner als vormittags. Der mittlere Unterschied betrug 39 Skalenteile, also etwas mehr als den 600\textsuperscript{sten} Teil der gesamten Intensität. An den meisten dieser Tage wurde der Stand des Apparats auch vormittags um 9 Uhr aufgezeichnet. Unter 28 Tagen waren 23, an denen die Intensität um diese Stunde noch größer war als eine Stunde später. Nur an 5 Tagen war es umgekehrt: Der mittlere Unterschied betrug hier aber nur $11\frac{1}{2}$ Skalenteile, also etwas mehr als den 2000\textsuperscript{sten} Teil der gesamten Intensität.

\expandafter\def\csname HistoricalNoteText371\endcsname{Auf ähnliche Weise und mit demselben Ergebnis wurde später auch beim Beobachtungstermin vom 13.---14. November mit beiden Apparaten beobachtet.}
Zweitens bestätigen mehrere sehr umfangreiche Beobachtungsreihen, dass unregelmäßige, bisweilen sehr beträchtliche und in kurzen Zeitabständen wechselnde Störungen bei der Intensität nicht weniger häufig vorkommen als bei der Deklination, woran man freilich schon aufgrund der Analogie nicht zweifeln konnte. Dreimal wurden bereits über beträchtliche Zeit hinweg mit diesem Intensitätsapparat und gleichzeitig mit dem Magnetometer des magnetischen Observatoriums ununterbrochen Aufzeichnungen gemacht. Das geschah am 15. Juli von morgens 6 Uhr bis abends 6 Uhr, dann beim ordentlichen magnetischen Beobachtungstermin vom 29.---30. Juli und schließlich beim außerordentlichen Termin vom 31. August bis zum 1. September, in den letzten beiden Fällen jeweils 24 Stunden lang. Die Aufzeichnungen erfolgten stets von 5 zu 5 Minuten. Grafische Darstellungen der beiden Termine, in denen Kurven für die Änderungen sowohl der Intensität als auch der Deklination gezeichnet sind, machen dies deutlich. Die beiderseitigen Bewegungen haben zwar, wie sich von selbst versteht, nicht die geringste Ähnlichkeit miteinander. Es ist aber sehr auffällig, dass bei starken Störungen der Deklination meistens auch starke Störungen der Intensität auftraten\HistoricalNote{371}.
\expandafter\def\csname HistoricalNoteText376\endcsname{Die erste öffentliche Erwähnung dieser Versuche findet sich in den Gött. gel. Anz. 1834, Aug. 9. Vgl. \textsc{Schumachers} Jahrbuch für 1836, [S. 339 dieses Bandes].}
\expandafter\def\csname HistoricalNoteText376b\endcsname{Gött. gel. Anz. 1835, März 7. \textsc{Schumachers} Jahrbuch für 1836, [S. 341 dieses Bandes.]}
\par
\SourceLeaf{372}Die Darstellung der Änderungen von Deklination und Intensität in zwei getrennten Kurven gibt übrigens ein bei Weitem weniger anschauliches Bild vom Verlauf der Störungen, als wenn man beide in einer einzigen Kurve zusammenführt. Worauf es dabei ankommt, erkennt man am klarsten auf folgende Weise. Die erdmagnetische Kraft zu einem beliebigen Zeitpunkt lässt sich vollständig durch eine einzige gerade Linie darstellen. Gemeint ist dabei wie immer stillschweigend ihr horizontaler Anteil. Die Länge der Linie ist proportional zur Intensität, und sie bildet mit einer festen Geraden einen Winkel, der der Deklination gleich ist. Will man die Kraft zu mehreren aufeinanderfolgenden Zeitpunkten darstellen, lässt man den Anfangspunkt der verschiedenen Geraden unverändert. Dann kommen nur ihre Endpunkte in Betracht, die man mit den entsprechenden Zeiten bezeichnen und durch eine Linie verbinden kann. Die geraden Radien selbst zeichnet man überhaupt nicht ein. Bei einem auch nur einigermaßen geeigneten Maßstab liegt sogar der gemeinsame Anfangspunkt stets weit außerhalb der Zeichnung. Dieses Vorgehen eröffnet uns zugleich eine neue Sicht auf solche Änderungen der beiden magnetischen Elemente. Sie sind in der Tat nur die beiden horizontalen Komponenten jener vergleichsweise immer sehr kleinen Störkraft, der die mittlere erdmagnetische Kraft in jedem Augenblick unterliegt. Dabei zerlegt man die Störkraft in zwei Richtungen, die eine \textit{im} magnetischen Meridian, die andere senkrecht dazu. Die zweite Komponente liefert unmittelbar das Magnetometer, die erste der neue Apparat. Vor dem Zeichnen müssen lediglich beide auf einen gemeinsamen Maßstab gebracht werden.

Bei der Anwendung dieser an sich so anschaulichen Darstellungsweise muss hier noch ein Umstand erwähnt werden. Wenn häufige, schnell wechselnde Störungen auftreten, lässt sich der Verlauf eines ganzen Tages kaum in einer einzigen Zeichnung übersichtlich darstellen, weil die Kurve dann sehr viele Verschlingungen aufweist. In diesem Fall muss man kürzere Zeitabschnitte einzeln zeichnen.

Vergleichen wir, was der neue Apparat und das Magnetometer leisten, zeigt sich, dass sie sich für \textit{einige} Zwecke gegenseitig ergänzen müssen, während sie sich für andere gleichermaßen einsetzen lassen. Zur Bestimmung der absoluten Deklination eignet sich nur das Magneto\SourceLeaf{373}meter. Die Änderungen der Deklination, insbesondere die schnell wechselnden, lassen sich mit beiden Apparaten verfolgen. Zur Bestimmung der absoluten Intensität eignen sich beide Apparate, obwohl die Anwendung des Magnetometers etwas weniger aufwendig ist, als es der alleinige Einsatz des neuen Apparats wäre. Das Magnetometer für sich genommen kann aber nur den Mittelwert der Intensität über einen bestimmten Zeitraum angeben. Schnell wechselnde Änderungen der Intensität entgehen ihm völlig, während der neue Apparat sie höchst zufriedenstellend nachweist. Für alle anderen Anwendungen sind beide gleich gut brauchbar, etwa um Magnetstäbe hinsichtlich ihrer magnetischen Stärke zu vergleichen oder in Verbindung mit einem Multiplikator für galvanometrische und telegrafische Zwecke. Für die beiden letzten Zwecke hat der neue Apparat sogar einen beträchtlichen Vorteil, weil man ihn, wie schon bemerkt, beliebig nahe an einen astatischen Zustand bringen kann.

Ein paar Beispiele für die Empfindlichkeit des Apparats als Galvanometer dürfen hier wohl angeführt werden. Der Multiplikator um den Magnetstab enthält 610 Windungen aus mit Seide umsponnenem Kupferdraht. Schon in ihm allein muss ein galvanischer Strom mehr als 6000 Fuß Draht durchlaufen. Diese Länge steigt auf 13000 Fuß, wenn der Strom zugleich zum Physikalischen Kabinett fließt. Gewöhnlich werden aber noch andere Apparate in den Stromkreis eingebracht, sodass bei manchen Versuchen die gesamte Drahtlänge 40000 Fuß oder fast zwei Meilen beträgt. Dabei ist noch zu beachten, dass bei Weitem der größte Teil dieses Drahts viel dünner ist. Soweit die Stromstärke von dieser Länge abhängt, entspricht sie etwa 8 Meilen eines Drahts von derselben Dicke wie der Verbindungsdraht zwischen der Sternwarte und dem Physikalischen Kabinett. Trotz des so langen Stromkreises bewirken selbst die schwächsten galvanischen Kräfte am schweren Magnetstab eine Auslenkung, die sich wahrnehmen lässt und für genaue Messungen ausreicht. Das gilt zum Beispiel für den Thermogalvanismus, von dem manche Physiker irrtümlich annehmen, er könne einen sehr langen Stromkreis nicht durchdringen. Bei den hiesigen Vorrichtungen und mit einem thermogalvanischen Apparat besonderer Bauart reicht es aus, die Verbindungsstelle mit dem Finger zu berühren, um diese Wirkung hervorzubringen.

\SourceLeaf{374}Die Anwendung des Apparats auf die gewöhnliche Reibungselektrizität gibt Anlass zu einer weiteren interessanten Beobachtung. Bekanntlich hat \textsc{Colladon} entdeckt, dass sie, durch einen Multiplikator geleitet, die Magnetnadel ganz ähnlich auslenkt wie ein hydrogalvanisch erregter Strom. Seine Versuche wurden anfangs bezweifelt und später von \textsc{Faraday} bestätigt. Dieser Physiker hat als Erster gezeigt, dass selbst in einer sehr starken elektrischen Batterie nicht mehr Elektrizität entwickelt wird, als schon sehr schwache hydrogalvanische Erregungsmittel in wenigen Sekunden durch einen Leitungsdraht von mäßiger Länge treiben. Mit den hiesigen Apparaten war zwar ebenfalls schon vor mehreren Jahren durch Versuche bestätigt worden, dass die elektromagnetische Wirkung der Maschinenelektrizität tatsächlich vorhanden und gering ist. Es schien jedoch der Mühe wert, diese Versuche mit dem neuen, so viel empfindlicheren Apparat zu wiederholen. Statt eine Leidener Flasche oder eine Batterie solcher Flaschen durch den Drahtstromkreis zu entladen (wie es \textsc{Colladon} und \textsc{Faraday} getan hatten), verband man nur den Konduktor und das Reibkissen einer Elektrisiermaschine im Physikalischen Kabinett mit den Enden des zur Sternwarte führenden Drahtstromkreises. Zusammen mit dem Multiplikator war er 13000 Fuß lang. Die Elektrisiermaschine wurde anhaltend mit gleichförmiger Geschwindigkeit gedreht. Bei einer Umdrehung pro Sekunde hielt sie dadurch den fünfundzwanzig Pfund schweren Magnetstab im neuen Apparat in der Sternwarte bei einer Auslenkung von 144 Skalenteilen (etwas über 50 Winkelminuten). Diese war je nach der Richtung, in der die Elektrizität den Multiplikator durchströmte, positiv oder negativ. Die Versuche verliefen mit aller nur wünschenswerten Regelmäßigkeit. Besonders bemerkenswert erscheint dabei, dass die elektromagnetische Wirkung auch dann dieselbe blieb, wenn man den Stromkreis durch Einfügen weiterer Apparate auf eine ganze Meile verlängerte. Das könnte als wesentlicher Unterschied zu anderen, hydrogalvanisch, thermogalvanisch oder durch Induktion erzeugten Strömen erscheinen. Ihre Intensität, die sich in der Größe der elektromagnetischen Wirkungen zeigt, nimmt stets ab, je weiter man die Leitung verlängert. Ich finde darin aber lediglich eine schlagende Bestätigung der Theorie, nach der die ungleiche Intensität zweier galvanischer Ströme, die sich an ihrer ungleichen elektromagnetischen Wirkung zeigt, allein darin besteht, dass in einer bestimmten Zeit ungleiche Mengen Elektrizität durch jeden Querschnitt der Leitung strömen.
\SourceLeaf{375}Bei den anderen Erzeugungsarten entwickelt eine gegebene elektromotorische Kraft in gegebener Zeit umso weniger Elektrizität, je größer der Widerstand ist, den der längere Stromkreis dem Strom entgegensetzt. Bei unserem Versuch hingegen hängt die Menge der bewegten Elektrizität allein vom Betrieb der Maschine ab. \textit{Alle} Elektrizität, die in Form von Funken auf den Konduktor überspringt, muss den ganzen Stromkreis durchlaufen, ob er kurz oder lang ist, um sich mit der entgegengesetzten Elektrizität des Reibkissens auszugleichen.

Um auch den Vorzug des neuen Apparats vor dem Magnetometer bei der elektromagnetischen Telegrafie zeigen zu können, ist zunächst etwas genauer zu betrachten, \textit{wie} telegrafische Zeichen durch galvanische Ströme hervorgebracht werden.

\expandafter\def\csname HistoricalNoteText375\endcsname{Nach einer Mitteilung, die mir Herr von \textsc{Humboldt} zukommen ließ, hatte \textsc{Bétancourt} schon zehn Jahre früher einen Drahtstromkreis von Aranjuez nach Madrid gezogen. Darin sollte die Entladung einer Leidener Flasche zur telegrafischen Übermittlung eines Zeichens dienen. Obwohl nähere Umstände über den Erfolg nicht bekannt zu sein scheinen, kann man am Gelingen eines solchen Versuchs bei zweckmäßiger Ausführung nicht zweifeln. Eine solche Methode müsste sich aber wohl stets darauf beschränken, auf eine oder ein paar vorher vereinbarte Fragen ein Ja oder Nein zu signalisieren.}
Sobald bekannt war, dass sich die Wirkungen einer \textsc{Volta}schen Säule durch einen sehr langen Stromkreis fortpflanzen, lag der Gedanke nahe, diese Naturkräfte zur Telegrafie zu nutzen. Schon vor fast 30 Jahren\HistoricalNote{375}, als man erst einen kleinen Teil der galvanischen Wirkungen kannte, schlug \textsc{Sömmerring} vor, dafür die Gasentwicklung zu nutzen. Für zusammengesetzte Signale eignen sich die erst später bekannt gewordenen magnetischen Wirkungen galvanischer Ströme weitaus besser. Es ist jedoch auffällig, dass nach \textsc{Oersted}s Entdeckung etliche Jahre verstrichen, bevor jemand an diese Verwendung gedacht zu haben scheint. Ein gründliches Urteil über die Anwendbarkeit im großen Maßstab setzt freilich genaue quantitative Kenntnisse darüber voraus, wie die Länge und die Beschaffenheit der Leitungsdrähte galvanische Ströme schwächen. Vor \textsc{Ohm} und \textsc{Fechner} hatte man davon sehr unvollständige und unrichtige Vorstellungen. Im Jahr 1833 wurde zwischen der hiesigen Sternwarte und dem Physikalischen Kabinett eine Drahtverbindung hergestellt. Dies geschah vor allem, um ähnliche Untersuchungen über das Gesetz der Stärke galvanischer Ströme unter verschiedenen Bedingungen in großem Maßstab anstellen zu können. Das Ver\SourceLeaf{376}dienst der sehr schwierigen Ausführung dieser groß angelegten Einrichtung gebührt allein Herrn Professor \textsc{Weber}. Dieser Stromkreis wurde von Anfang an oft für telegrafische Zeichen verwendet, für einfache Zeichen zum täglichen Vergleich der Uhren und versuchsweise auch für zusammengesetzte Zeichen. Die Möglichkeit, Buchstaben, Wörter und ganze Sätze zu signalisieren. wurde dadurch schon damals zu einer offenkundigen Tatsache\HistoricalNote{376}. Bei diesen Versuchen verwendete man einen hydrogalvanischen Strom, der nur mit schwachen Mitteln erregt wurde, nämlich mit einem einfachen oder doppelten Plattenpaar und Wasser ohne Säurezusatz. Ich halte mich hier jedoch nicht mit einer ausführlichen Beschreibung des damals verwendeten Verfahrens auf, da ich es später durch ein ganz anderes ersetzt habe. Jenes Verfahren hatte weiterhin den Nachteil, dass sich mit unserem einfachen Stromkreis und der damaligen Einrichtung der Apparate, bei denen solche Versuche nur eine Nebensache waren, in einer Minute höchstens zwei Buchstaben signalisieren ließen. Auch mit einer abgeänderten, eigens für die Telegrafie bestimmten Einrichtung hätte sich diese Geschwindigkeit kaum steigern lassen, solange man nur einen einfachen Stromkreis benutzte. Die Länge des Stromkreises oder der Abstand seiner Endpunkte hat damit übrigens offensichtlich nichts zu tun. Mit einem mehrfachen Stromkreis wäre dagegen eine erhebliche Steigerung möglich gewesen. Es gab hier aber keinen hinreichenden Grund, einen solchen einzurichten. Zum einen konnte der Erfolg selbst gar nicht zweifelhaft sein, zum anderen hätte der eigentliche wissenschaftliche Nutzen eines solchen mehrfachen Stromkreises in keinem Verhältnis zu den beträchtlichen Kosten gestanden.

Die Theorie der Induktionsgesetze hat mich dagegen auf ein ganz anderes Verfahren geführt. Danach dient ein einfacher Stromkreis schon seit mehr als zwei Jahren mit vollkommenem Erfolg zu einer viel schnelleren Telegrafie. Ich darf wohl umso eher noch etwas dabei verweilen, als ich bisher keine Einzelheiten darüber öffentlich bekannt gemacht habe.

Die Vorrichtung, die ich Induktor nenne, habe ich schon vor mehreren Jahren an anderer Stelle beschrieben{\renewcommand{\thefootnote}{**)}\HistoricalNote{376b}}. Ich muss jedoch darauf hinweisen, dass anstelle des in der ersten Mitteilung beschriebenen Induktors mit 1050 Windun\SourceLeaf{377}gen und des später auf 3537 Windungen verstärkten inzwischen einer mit 7000 Windungen verwendet wird, in dem die Drahtlänge allein mehr als 7000 Fuß beträgt. Durch einen äußerst einfachen Handgriff mit diesem Induktor gehen zwei starke, entgegengesetzte galvanische Ströme rasch nacheinander durch den Leitungsdraht. Jeder dauert nur äußerst kurze Zeit. Man zieht den Induktor rasch von dem doppelten Magnetstab ab, über den er zunächst geschoben ist, und bringt ihn sofort wieder in seine vorige Lage, ohne ihn umzudrehen. Befindet sich irgendwo im Stromkreis eine von einem Multiplikator umgebene Magnetnadel, so erteilen ihr die beiden Ströme für einen Augenblick eine sehr hohe Geschwindigkeit, die dann sofort wieder vollständig aufgehoben wird. Die Nadel macht also eine sehr schnelle, aber nur kleine Bewegung nach rechts oder links, wie man es wünscht, und steht gleich darauf wieder ganz still.

Es versteht sich von selbst, dass sich solche ruckartigen Bewegungen in verschiedener Folge kombinieren und zur Übermittlung von Buchstaben nutzen lassen. Um die Zeichen möglichst schnell und präzise zu geben und sie auf der anderen Seite leicht und schnell zu lesen, ist allerdings eine gewisse Übung nötig. Doch auch ohne besondere Übung lassen sich, wie wiederholte Erfahrungen gezeigt haben, in einer Minute ohne Weiteres etwa sieben Buchstaben signalisieren. Würde man für den Handgriff besondere \textit{mechanische} Vorrichtungen vorsehen, ließen sich Schnelligkeit und Präzision zweifellos noch erheblich steigern.

Gerade bei dieser Art des Telegrafierens hat der neue Apparat aus den folgenden Gründen einen bedeutenden Vorzug vor dem Magnetometer. Die beiden entgegengesetzten Impulse, aus denen ein einfaches Zeichen besteht, sind ihrer \textit{Stärke} nach \textit{genau gleich}. Der zweite hebt daher genau so viel Geschwindigkeit auf, wie der erste erzeugt hat. Dennoch kann die Nadel zwischen den Zeichen nicht vollkommen in Ruhe sein, denn das ist nur in ihrer natürlichen Gleichgewichtsstellung möglich. Selbst wenn sie sich \textit{vor} einem Zeichen in dieser Stellung befindet, wird sie eben \textit{durch} das Zeichen ein wenig daraus ausgelenkt, wenn auch nur geringfügig. Die auf die Nadel wirkende Richtkraft sucht sie dann in diese Stellung zurückzuführen. Auch wenn ein einzelnes Zeichen so nur eine äußerst \SourceLeaf{378}schwache Bewegung hervorrufen kann, kann sich durch die Anhäufung sehr vieler Zeichen eine beträchtliche Entfernung von der natürlichen Gleichgewichtsstellung ergeben. Dadurch kann zwischen den Zeichen so viel Bewegung entstehen, dass die Zeichen etwas an Deutlichkeit verlieren. Wie man bei einigem Nachdenken leicht einsieht, wirkt sich diese Störung unter sonst gleichen Umständen ungünstiger aus, wenn die Nadel, an deren ruckartigen Bewegungen man die Zeichen erkennt, eine kurze Schwingungsdauer hat, als wenn sie eine lange hat. Die Störung macht sich daher am Magnetometer des magnetischen Observatoriums stärker bemerkbar als an dem früher in der Sternwarte aufgehängten Magnetometer mit der fünfundzwanzig Pfund schweren Nadel. Sie macht sich noch weniger am neuen Apparat bemerkbar, der jetzt dessen Stelle einnimmt, wenn sein Magnetstab in umgekehrter Lage als nahezu astatische Nadel eingerichtet ist. Selbst nach einer beträchtlichen Entfernung aus der Gleichgewichtsstellung versetzt die vergleichsweise schwache Richtkraft, die die Nadel dorthin zurücktreibt, sie tatsächlich in keine Bewegung, die die Zeichen erheblich stören würde. Zugleich wirkt der Strom im Multiplikator genauso stark auf sie und ruft daher ebenso große ruckartige Bewegungen hervor, als gehörte sie zu einem gewöhnlichen Magnetometer.

Gegen die Nachteile und Unannehmlichkeiten unerwünschter Schwingungen, sowohl bei dieser Art des Telegrafierens als auch bei manch anderer Anwendung magnetischer Apparate, leistet eine besondere Vorrichtung, die ich vor Kurzem habe anfertigen lassen, äußerst nützliche Dienste. Ich nenne diese Vorrichtung einen \textit{Dämpfer}, da sie Schwingungen, die sonst viele Stunden andauern und dabei nur sehr langsam abnehmen würden, in sehr kurzer Zeit vollständig zum Stillstand bringt. Der zunächst nur für das Magnetometer des magnetischen Observatoriums angefertigte Dämpfer erfüllt diese Aufgabe in so hohem Maß, dass die größten Schwingungen in wenigen Minuten ganz erlöschen. Eine ähnliche Vorrichtung lässt sich bei jeder schwingenden Nadel anbringen, bei einem Magnetometer ebenso wie bei dem neuen Apparat, von dem hier die Rede ist. Bei allen Apparaten, die ernsthaft zum Telegrafieren nach dem hier beschriebenen Verfahren eingesetzt werden sollen, wird sie ein wesentlicher Bestandteil sein müssen. Eine ausführlichere Erklärung dieser Vorrichtung würde jedoch zu weit vom gegenwärtigen Gegenstand wegführen.

\SourceLeaf{379}Der neue Apparat hat bisher keine besondere Bezeichnung erhalten. Nach seiner wichtigsten Anwendung könnte man ihn Intensitätsmesser nennen. Da er aber zu ebenso vielfältigen und genauen magnetischen Messungen dient wie das Magnetometer, könnte er wohl ebenso gut \textit{dieselbe} Bezeichnung beanspruchen. Der wesentlichste Unterschied besteht darin, dass der neue Apparat an \textit{zwei} Fäden aufgehängt ist. Dadurch erhält man eine \textit{neue} Richtkraft, mit der die magnetische Richtkraft kommensurabel ist. Die übrigen Unterschiede, insbesondere die Art der Spiegelbefestigung und die Mittel zur Messung der Drehung der einzelnen Teile gegeneinander, sind notwendige Bedingungen, die sich aus den angestrebten Zwecken von selbst ergeben. Man könnte den neuen Apparat daher ein \textit{Bifilar-} oder \textit{Bipensil-Magnetometer} nennen, um ihn vom älteren Instrument, dem einfachen oder Unifilar-Magnetometer, zu unterscheiden. Ich darf wohl meine Überzeugung aussprechen, dass seiner weiteren Verbreitung und besonders seinem Einsatz bei den Terminbeobachtungen neben dem einfachen Magnetometer an mehreren weit voneinander entfernten Orten bald bedeutende Fortschritte in unserer Kenntnis der wunderbaren Störungen des Erdmagnetismus folgen werden.

